\end{ejercicio}

\begin{solucion}
\[
\varepsilon_{xx}=\frac{\partial u}{\partial x}=\boxed{2\times10^{-3}}.
\]
\[
\varepsilon_{yy}=\frac{\partial v}{\partial y}=\boxed{1\times10^{-3}}.
\]
Además:
\[
\varepsilon_{xy}
=\frac12\left(\frac{\partial u}{\partial y}+\frac{\partial v}{\partial x}\right).
\]
Por tanto:
\[
\varepsilon_{xy}
=\frac12\left(10^{-3}-0.5\times10^{-3}\right)
=\boxed{2.5\times10^{-4}}.
\]
La deformación angular ingenieril es:
\[
\gamma_{xy}=2\varepsilon_{xy}=\boxed{5\times10^{-4}}.
\]
\end{solucion}

\begin{ejercicio}[6: Viga con carga uniforme y fuerza puntual]
Una viga simplemente apoyada de longitud $L=\qty{6}{\metre}$ está sometida a $q=\qty{2}{\kilo\newton\per\metre}$ en toda su longitud y a una fuerza $P=\qty{6}{\kilo\newton}$ en el centro.

Obtener reacciones, cortante y momento flector.
\end{ejercicio}

\begin{solucion}
El problema es simétrico.

Carga total:
\[
qL+P=12+6=\qty{18}{\kilo\newton}.
\]
Por simetría:
\[
\boxed{R_A=R_B=\qty{9}{\kilo\newton}}.
\]
Para $0\leq x<3$:
\[
\boxed{V(x)=9-2x},
\qquad
\boxed{M(x)=9x-x^2}.
\]
Justo antes de la fuerza puntual:
\[
V(3^-)=3.
\]
La carga $P$ produce un salto:
\[
V(3^+)=3-6=-3.
\]
Para $3<x\leq6$:
\[
\boxed{V(x)=3-2x}.
\]
El momento es continuo:
\[
M(x)=9x-x^2-6(x-3),
\]
es decir:
\[
\boxed{M(x)=3x-x^2+18}.
\]
Como el cortante cambia de positivo a negativo en $x=3$:
\[
M_{\max}=M(3)=\boxed{\qty{18}{\kilo\newton\metre}}.
\]

\begin{center}
\begin{minipage}{0.92\linewidth}
\centering
\begin{tikzpicture}[x=1.0cm,y=0.95cm,>=Latex,
  loadarrow/.style={-{Latex[length=2.2mm]},BookOrange,line width=.85pt},
  reactarrow/.style={-{Latex[length=2.2mm]},BookBlue,line width=.95pt}
]
  % Viga
  \draw[BookGray,line width=2.0pt] (0,0)--(8,0);

  % Apoyos
  \draw[BookGray,line width=.85pt] (0.55,0)--(0.18,-0.52)--(0.92,-0.52)--cycle;
  \draw[BookGray,line width=.85pt] (7.45,0)--(7.08,-0.46)--(7.82,-0.46)--cycle;
  \draw[BookGray,line width=.75pt] (7.18,-0.64) circle (0.09);
  \draw[BookGray,line width=.75pt] (7.72,-0.64) circle (0.09);
  \draw[BookGray,line width=.75pt] (6.98,-0.79)--(7.92,-0.79);

  % Carga distribuida uniforme
  \draw[BookOrange,line width=1.0pt] (0.70,1.30)--(7.30,1.30);
  \foreach \x in {0.70,1.20,...,7.20}{
    \draw[loadarrow] (\x,1.30)--(\x,0.10);
  }
  \node[font=\sffamily\small,text=BookOrange!90!black] at (2.20,1.62)
    {$q=\qty{2}{\kilo\newton\per\metre}$};

  % Carga puntual
  \draw[loadarrow,line width=1.15pt] (4.00,2.10)--(4.00,0.10);
  \node[font=\sffamily\small,text=BookOrange!90!black] at (4.00,2.38)
    {$P=\qty{6}{\kilo\newton}$};

  % Reacciones
  \draw[reactarrow] (0.55,-1.38)--(0.55,-0.08);
  \draw[reactarrow] (7.45,-1.38)--(7.45,-0.08);
  \node[font=\small,text=BookBlue,left] at (0.42,-1.25) {$R_A$};
  \node[font=\small,text=BookBlue,right] at (7.58,-1.25) {$R_B$};

  % Cotas
  \draw[BookGray!75,|<->|,line width=.55pt] (0.55,-1.75)--(4.00,-1.75)
    node[midway,fill=white,inner sep=1.5pt,font=\scriptsize] {\qty{3}{\metre}};
  \draw[BookGray!75,|<->|,line width=.55pt] (4.00,-1.75)--(7.45,-1.75)
    node[midway,fill=white,inner sep=1.5pt,font=\scriptsize] {\qty{3}{\metre}};

  \node[font=\small] at (0.55,-0.72) {$A$};
  \node[font=\small] at (7.45,-0.72) {$B$};
\end{tikzpicture}
\captionof{figure}{Viga simplemente apoyada del ejercicio 6, con carga uniforme y fuerza puntual centrada.}
